Celem niniejszego rozdziału jest zbadanie tworzenia tożsamych programów obiektowych o różnych współczynnikach sprzężenia.
Takie narzędzie jest konieczne do kontroli tej zmiennej niezależnej eksperymentu.

W tym rozdziale zostanie przedstawiona definicja tożsamych programów w rozumieniu klasycznej semantyki operacyjnej,
a następnie jej dostosowanie do kontekstu programowania obiektowego.
Przybliżony zostanie sposób redukcji siły sprzężenia pomiędzy klasami za pomocą powszechnie uznanych metod
refaktoryzacji kodu źródłowego.
Następnie w drodze badania możliwości manipulacji zmienną niezależną zostanie przedstawiona metoda celowego zwiększenia
przypadkowej złożoności: odwróconej refaktoryzacji.
Przedstawione zostaną przykłady takiego zabiegu porównane za pomocą miar sprzężenia wspomnianych w rozdziale~\ref{cha:coupling}.
Na zakończenie rozdziału zostanie omówiony przykład programu \texttt{JavaBus}, użytego do eksperymentu, którego
metodę badawczą przedstawiono w rozdziale~\ref{cha:metoda}, a wyniki w rozdziale~\ref{cha:wyniki}.

\section{Pojęcie tożsamości programów}

Jednym z głównych zadań postawionych przed pracą jest utworzenie dwóch alternatywnych wersji programu o różnych
stopniach sprzężenia pomiędzy klasami.
Utworzenie alternatywy nie może być przeprowadzone przez dowolną modyfikację programu wyjściowego.
Zmiany kodu źródłowego programu mogą zmienić program tak, że nie spełnia on swojej założonej funkcji.
Jeżeli do tego dojdzie, porównywanie zjawisk występujących przy wprowadzaniu zmian będzie nieuzasadnione.
Każde odstępstwo od wspólnego punktu odniesienia niesie ze sobą dodatkowe zniekształcenia w przeprowadzeniu
kontrolowanego eksperymentu, a ich znaczenie omówiono szczegółowo w rozdziale~\ref{cha:metoda}.

\begin{figure}[h]
    \centering
    \begin{minipage}[t]{0.45\linewidth}
        \lstset{
            basicstyle=\ttfamily\small,
            frame=single,
            numbers=left,
            numberstyle=\tiny,
            language=Java
        }
        \begin{lstlisting}
String rev(String s) {
 String r="";
 for(int i=s.length()-1;i>=0;i--){
  r+=s.charAt(i);
 }
 return r;
}
        \end{lstlisting}
    \end{minipage}
    \hfill
    \begin{minipage}[t]{0.45\linewidth}
        \lstset{
            basicstyle=\ttfamily\small,
            frame=single,
            numbers=left,
            numberstyle=\tiny,
            language=Java
        }
        \begin{lstlisting}
String rev(String s) {
 String reversed="";
 for(int i=s.length()-1;i>=0;i--){
  r = addChar(r, s.charAt(i);
 }
 return reversed;
}
String addChar(String r, char c){
 return r + c;
}
        \end{lstlisting}
    \end{minipage}
    \caption{Dwie funkcje \texttt{rev}, które mogą być rozważane jako tożsame programy}
    \label{fig:podobny_kod}
\end{figure}

Z powyższego formułuje się następująca intuicja: należy tak zmodyfikować program, aby zmiany nie wpłynęły na funkcję
programu, jednocześnie w pożądany sposób zmieniły jego strukturę.
Przykładem czynności nienaruszających funkcji programu mogą być zmiana nazw zmiennych używanych w metodach lub
dodanie prywatnej metody w celu wydzielenia bloku z istniejącego kodu.
Przykład na Rysunku~\ref{fig:podobny_kod} ilustruje takie modyfikacje, programista naturalnie stwierdzi, że oba
programy są takie same pomimo widocznych różnic.
Takie czynności zmiany kodu, powszechnie znane jako refaktoryzacje (szczegółowo omówione w dalszych podrozdziałach),
nie wpływają na przebieg, a tym samym na wynik programu.

\begin{figure}[h]
    \centering
    \begin{minipage}[t]{0.45\linewidth}
        \lstset{
            basicstyle=\ttfamily\small,
            frame=single,
            numbers=left,
            numberstyle=\tiny,
            language=Java
        }
        \begin{lstlisting}
String rev(String s) {
 String r = "";
 for(int i=s.length()-1;i>=0;i--){
  r = s.charAt(i) + r;
 }
 return r;
}
        \end{lstlisting}
    \end{minipage}
    \hfill
    \begin{minipage}[t]{0.45\linewidth}
        \lstset{
            basicstyle=\ttfamily\small,
            frame=single,
            numbers=left,
            numberstyle=\tiny,
            language=Java
        }
        \begin{lstlisting}
String rev(String s) {
 String r = "";
 for(int i=s.length()-1;i>=0;i--){
  r = r + s.charAt(i);
 }
 return r;
}
        \end{lstlisting}
    \end{minipage}
    \caption{Dwie bardzo podobne do siebie funkcje \texttt{rev}, które na pewno nie są tożsame}
    \label{fig:rozny_kod}
\end{figure}

Inny przypadek brzegowy: zmiana kolejności wynikowych elementów w programie sortującym sprawia, że program zdecydowanie
pełni inną funkcję i nie może być rozważany jako alternatywna struktura kodu rozwiązująca ten sam problem.
Innym przykładem jest Rysunek~\ref{fig:rozny_kod}, gdzie przedstawiony kod źródłowy jest bardzo podobny, ale
rezultat jest różny, zatem te funkcje nie mogą być tożsamymi programami.
W wielu algorytmach sortowania takie przekształcenie programu polega na zmianie znaku pojedynczej operacji porównania.
W związku z powyższym rozmiar modyfikacji nie świadczy o znaczeniu zmiany.

Wynika stąd potrzeba określenia narzędzia badawczego, które będzie w stanie stwierdzić, że oba programy, pomimo różnic
strukturalnych, są takie same.
Nieodzownym jest wbudowanie tego narzędzia do eksperymentu w niniejszej pracy jako kryterium przygotowanej pary
programów.

Rozumowaniu ,,takich samych dwóch programów'' podanemu powyżej brakuje ścisłej definicji.
Choć pojawia się pojęcie ,,zmiany kodu źródłowego tak, by nie zmienić programu'', ciężko je uznać za weryfikowalne.
Sugeruje to, że programy mogą być ,,podobne'' lub nawet ,,takie same'' pomimo różnic syntaktycznych.
W tym celu autor pracy wprowadza poniżej definicję programów tożsamych, w pierwszej kolejności przedstawiając
definicje występujące w literaturze.

\subsection{Definicje tożsamości w literaturze}

Tożsamość programów jest pojęciem wywodzącym się z informatyki teoretycznej, dokładniej z teorii automatów.
Jest to zagadnienie nietrywialne, stosowane do analizy programów pod kątem własności stopu, optymalizacji i
weryfikacji poprawności~\cite{tzevelekos2012program}.
Literatura skupia się na teoretycznym zastosowaniu do analizy programów napisanych w formalnych językach
programowania, takich jak \texttt{nu-calculus}.
Są to zastosowania nieadekwatne bezpośrednio do programów napisanych w językach programowania wysokiego poziomu,
niemniej na ich podstawie można przeprowadzić syntezę definicji osadzonej w wymaganym kontekście.

Artykuł~\cite{tzevelekos2012program} wprowadza pojęcie tożamości programów jako ,,dwa programy są
tożsame obserwacyjnie, lub po prostu tożsame, jeżeli można ich używać zamiennie użyte w dowolnym kontekście wykonania
programu bez wpływania na obserwowalne zachowanie tego kontekstu''.
Kontekstem wykonania jest pamięć, czyli stan programu, a także wejście-wyjście środowiska, w którym program jest
uruchomiony.
Do tej definicji wymagane są dodatkowe założenia w kontekście środowiska obsługującego błędy, więc aby programy były
tożsame, musi zajść jeden z warunków~\cite{churchill2019semantic}:
\begin{itemize}
    \item oba programy kończą się normalnie z podobnymi stanami stosów i rejestrów wyjścia, lub
    \item każdy z programów kończy się błędem wykonania, lub zamiast tego wykonuje się w nieskończonej pętli.
\end{itemize}

W tej definicji oba tożsame programy muszą zakończyć się w tym samym stanie albo ulec awarii.
Formalne języki programowania, takie jak \texttt{nu-calculus} rozważany w przytaczanej pracy, posiadają bardzo
ścisłą definicję stanu, który wymaga uogólnienia na inne języki programowania.
Warunek awarii nie odpowiada zastosowaniu w niniejszej pracy --- zakłada, że oba programy mogą pełnić w założeniu
inne funkcje, ale jeżeli kończą się błędem lub nie spełniają warunku stopu, to są równoważne.
Przy takim założeniu można porównywać ze sobą dowolne programy, o ile zawierają błędy prowadzące do powyższych.

Artykuł poświęcony \cite{plotkin1981structural} definiuje tożsamość programów \(e\) i \(e'\) dla stanu \(\sigma\) jako:
\begin{equation}
    \label{eq:equivalence}
    \forall\sigma.eval(e, \sigma) = eval(e', \sigma)
\end{equation}

To oznacza, że oba tożsame programy dla dowolnego wejścia pozostawiają ten sam stan.
Problemem w złożonych nowoczesnych programach jest ustalenie i porównywanie tego stanu:
może na niego składać się wiele elementów niezwiązanych z wykonaniem samego programu, zaczynając od niederministycznego
przełączania procesów w systemie operacyjnym, na nieprzewidywalności alokacji stron pamięci kończąc.
Stan musi zatem zostać dokładnie określony, pomijając elementy, które nie mają znaczenia przy analizie pod kątem
tożsamości programów.
Podobne rozumienie testowania tożsamości występuje w innych publikacjach, dekomponując stan na dane wejściowe i
wyjściowe~\cite{corbett2000bandera, tzevelekos2012program}.

Idea tożsamości programów dla wszystkich definicji jest spójna i oznacza równoważność wyniku obu programów na wejście.
W przypadku poniższej pracy należy osadzić powyższe w kontekście programów napisanych w językach wysokiego poziomu,
biorąc pod uwagę ograniczenia i narzędzia dostępne w tym środowisku.

\subsection{Definicja tożsamości programu obiektowego}

Cytowany już~\cite{tzevelekos2012program} wspomina również o stwierdzaniu tożsamości dwóch programów na podstawie
zbioru wymagań, który jest weryfikowany dla dwóch programów, wobec czego stwierdza się tożsamość zachowania.
Ta definicja jest bliższa potrzebom niniejszej pracy, którą można sformułować jako:

\begin{definition}
    \label{def:equivalence}
    Program \(A\) jest tożsamy z \(B\), kiedy \(A\) spełnia wymagania \(B\) i \(B\) spełnia wymagania \(A\).
\end{definition}

Do definicji~\ref{def:equivalence} potrzeba jeszcze rozwinąć pojęcie \textit{spełnienia wymagań}:

\begin{definition}
    \label{def:requirements}
    Program \(A\) spełnia wymagania \(B\) kiedy wszystkie wymagania wobec programu \(A\) zostają spełnione dla \(B\)
\end{definition}

Pojęcie \textit{spełnienia wymagań} nie jest precyzyjne.
Formułowanie wymagań wobec programów jest obszernym tematem, którym zajmuje się inżynieria wymagań oprogramowania~\cite{macaulay2012requirements}.
Na potrzeby tej pracy można przyjąć, że wymaganie wobec programu może przybrać formę określenia wejścia i
wyjścia programu, w zależności od jego kontekstu wykonania.
Oznacza to, że jeżeli program jest uruchamiany jako polecenie terminala, to wymaganiem będzie para komendy
wejściowej i oczekiwanego wyniku, a zbiór takich wymagań będzie ich zestawem.
Podobnie, jeżeli program wyposażony jest w interfejs użytkownika i ma działać jako taki, to wymaganiem będzie
scenariusz testowy, opisujący krok po kroku interakcje użytkownika z interfejsem i oczekiwane rezultaty takiej
interakcji.

\begin{lstlisting}[language=Java,
    caption={Przykładowy test JUnit dla metody statycznej \texttt{rev}},
    label={lst:sample_unit_test}]
@Test
void shouldReturnReversedString() {
    String input = "abc";

    String result = rev(input);

    assertEquals("cba", input);
}
\end{lstlisting}

Dla programów, które są bibliotekami, wymagania mogą manifestować się w formie testów jednostkowych: scenariuszy
użycia obiektów dostarczanych przez bibliotekę wraz z oczekiwanymi rezultatami.
Listing~\ref{lst:sample_unit_test} przedstawia test jednostkowy skonstruowany dla wcześniej przedstawionej metody
statycznej \texttt{rev}.
Ze względu na to, że ciało metody reprezentującej test może być dowolne, należy dodać wymagania wobec testu
jednostkowego precyzujące test jednostkowy, który może być rozważany jako wymaganie.
Test jednostkowy opisujący wymaganie wobec programu musi:
\begin{itemize}
    \item jednoznacznie precyzować wejście programu,
    \item wywoływać kod opisywanego programu oraz
    \item weryfikować dane wyjściowe pod kątem oczekiwanych rezultatów.
\end{itemize}

Niespełnienie przynajmniej jednego z przytoczonych wymagań skutkuje odrzuceniem takiego testu jako testu jednostkowego
opisującego wymaganie na program.
Nie jest to sytuacja, która dotyczy wyłącznie wadliwych testów.
Testy dymne (ang.\@ \textit{smoke tests}) stosowane są w ramach wykrywania wad, ale z reguły nie weryfikują wyjścia
programu, najczęściej odtwarzając główne scenariusze użycia pod kątem wystąpienia błędów krytycznych~\cite{chauhan2014smoke}.

Dane wejściowe i wyjściowe testu nie muszą być wyłącznie parametrami metod: może to być stan innych obiektów,
statycznych atrybutów, a także urządzeń wejścia/wyjścia.

Założenie, że testy jednostkowe opisują wszystkie wymagania na program, nie jest pozbawiony wad.
Przede wszystkim należy założyć, że zbiór wymagań na program jest kompletny.
Jest to założenie nieweryfikowalne, w szczególności, że dla programów o średniej złożoności prawie nigdy nie formułuje
się pełnych wymagań, gdyż część z nich wynika z kontekstu wykonania~\cite{jkuchta2016}.
Z tak założonego pełnego zbioru wymagań należy utworzyć \textbf{pełen} zestaw testów jednostkowych.
Innymi słowy, zestaw testów jednostkowych musi w pełni odwzorowywać wymagania na program.

Innym sposobem, by rozwiązać problem kompletności zestawu testów jednostkowych opisujących wymagania, jest zbadanie
zestawu pod kątem pokrycia kodu.
Każdy uruchomiony test jednostkowy pozostawia ślad wykonania kodu produkcyjnego (\textit{kodem produkcyjnym} nazywa
się kod realizujący wymaganie i będący przedmiotem właściwego programu), następnie środowisko uruchomieniowe testu
sumuje ślady i określa, ile z linii kodu zostało uruchomionych.
Ta prosta metryka, zwana też wskaźnikiem pokrycia linii kodu, pomaga programistom na ustalenie, który kod nie został
jeszcze wywołany.
Niestety, taki sposób pomiaru nie wskazuje, ile odgałęzień kodu zostało uruchomionych, a posługiwanie się tą metryką
pozwala na wykrycie jedynie 19\% błędów w przypadku 100\% pokrycia linii kodu~\cite{7272926}.
Aby prześledzić dokładniej uruchomione odgałęzienia, stosuje się metrykę pokrytych gałęzi kodu --- czyli określa się,
ile procent przypadków zostało uruchomionych.
W związku z powyższym dla zestawu testów spełniających warunki testów jednostkowych opisujących wymagania na program
należy wskazać minimalny obszar pokrycia gałęzi kodu.
W przypadku poniższej pracy można założyć, że zestaw testów opisujących kompletne wymagania na program pokrywają
100\%~gałęzi kodu produkcyjnego.

Podsumowując, wprowadza się następującą definicję tożsamości programów, które mogą zostać opisane testami
jednostkowymi w \texttt{Java}:
\begin{definition}
Program \(A\) jest tożsamy z programem \(B\) jeżeli istnieje wspólny zestaw testów jednostkowych opisujących wymagania
pokrywający 100\% gałęzi programów \(A\) i \(B\).
\end{definition}

\section{Redukcja sprzężenia klas za pomocą refaktoringu}

Program tożsamy z innym może składać się z innych klas wewnętrznych, lub metody jego mogą być poprzesuwane.

Refactoring Fowlera opowiada o takich modyfikacjach programu i sugeruje takie rozwiązania.

Jeżeli posiada się program, który realizuje założone wymagania, to można zastosować odwrócenie refactoringów Fowlera do
zwiększenia złożoności.

Oba repozytoria zostały pozbawione poprzedniej historii tak, aby badany nie mógł możliwości odkryć zmian, które zaszły w projekcie:

\section{Odwrócona refaktoryzacja jako sposób na modyfikację programu}

\section{Modyfikacja rogramu \texttt{JavaBus} metodą odwróconej refaktoryzacji}

\begin{itemize}
    \item https://github.com/puradawid/silver-winner (współczynnik sprzężenia klas: 1)
    \item https://github.com/puradawid/verbose-parakeet (współczynnik sprzężenia klas: 3)
\end{itemize}

\begin{lstlisting}[language=Java,
    caption={Klasa InMemoryBus o współczynniku sprzężenia z klasą DeliveryStrategy równym 1},
    label={lst:inmemorybus_1}]
package edu.pw.javabus;

public class InMemoryBus implements Bus {

    private final RegisteredConsumers registeredConsumers = new RegisteredConsumers();

    private final DeliveryStrategy deliveryStrategy;

    InMemoryBus(DeliveryStrategy deliveryStrategy) {
        this.deliveryStrategy = deliveryStrategy;
    }

    @Override
    public <T extends Message> void register(Consumer<T> consumer, Topic topic, Class<T> messageType) {
        registeredConsumers.add(consumer, topic, messageType);
    }

    @Override
    public <T extends Message> void unregister(Consumer<T> consumer, Topic topic) {
        registeredConsumers.remove(consumer, topic);
    }

    @Override
    public <T extends Message> void send(Topic topic, T message) {
        deliveryStrategy.deliver(registeredConsumers, topic, message);
    }

}
\end{lstlisting}

\begin{lstlisting}[language=Java,
    caption={Klasa InMemoryBus o współczynniku sprzężenia z klasą DeliveryStrategy równym 3},
    label={lst:inmemorybus_3}]
package edu.pw.javabus;

import java.util.Iterator;

public class InMemoryBus implements Bus {

    private final RegisteredConsumers registeredConsumers = new RegisteredConsumers();

    private final MessageDelivery messageDelivery;

    InMemoryBus(MessageDelivery messageDelivery) {
        this.messageDelivery = messageDelivery;
    }

    @Override
    public <T extends Message> void register(Consumer<T> consumer, Topic topic, Class<T> messageType) {
        registeredConsumers.add(consumer, topic, messageType);
    }

    @Override
    public <T extends Message> void unregister(Consumer<T> consumer, Topic topic) {
        registeredConsumers.remove(consumer, topic);
    }

    @Override
    public <T extends Message> void send(Topic topic, T message) {
        boolean isLastDeliveryASuccess;
        boolean anySuccess = false;
        while (topic != null) {
            Iterator<Consumer<? extends Message>> iterator = registeredConsumers.find(topic, message.getClass()).iterator();
            while (iterator.hasNext()) {
                Consumer<? extends Message> nextConsumer = iterator.next();
                if (messageDelivery.shouldCallThisConsumer(nextConsumer, message, anySuccess)) {
                    isLastDeliveryASuccess = new ConsumeMethodObject(nextConsumer).consume(message).isSuccess();
                } else {
                    continue;
                }
                if (messageDelivery.shouldStop(isLastDeliveryASuccess)) {
                    return;
                }
                anySuccess = anySuccess || isLastDeliveryASuccess;
            }
            topic = messageDelivery.shouldGetParentTopic(anySuccess) ? topic.getParent() : null;
        }
    }

}

\end{lstlisting}

